% New result after checkpoint04; independent review pending.
\subsection{A prescribed finite test can coexist with first failure at class three}

\begin{theorem}[Selected A.72 pairs first fail at integral class three]
\label{thm:ghost-a72-first-class3}
Let \(N\triangleleft\Pi\) be open and normal in the full
arithmetic group \(E_U\), and put \(Q=\Pi/N\).
There is an ordinary infinite-support completed pair
\((A,B)\) produced by the recursion of SECTIONS A.72,
with a continuous section
\[
 \tau:G_{\mathbb Q}\longrightarrow
 E_U/(N\cap[\Pi,\Pi])
   =E_Q\times_{E_{Q^{\rm ab}}}E_{\rm ab},
 \tag{AQ.1}
\]
whose \(Q\)-component is the section of three and whose
restriction at every finite place is geometrically
conjugate to the image of its ordinary local point.
At the real place it has the image of three.
Nevertheless the unrestricted full integral nilpotent
section fibre above \((A,B)\) is nonempty in class two
and empty in class three.

Consequently no one fixed finite geometric quotient,
even retained jointly with the full abelianization and
all the local conjugacies in (AQ.1), decides integral
class-three liftability for ordinary completed pairs.
\end{theorem}
\begin{proof}
We use the fully reconstructed recursion and component
maps of Theorem~\ref{thm:ghost-finite-quotient-survival},
and specify an additional choice at its first prime.
Retain its pointed \(Q\)-cover
\(Y\to U\), its rational point \(y_0\) above three,
and its finite Galois field \(F\).
That field defines all boundary points and tangent
leading coefficients required in (FQ.3), all geometric
deck transformations, the \(|Q|\)-th roots of three
and minus two, and the required roots of the leading
coefficients. Put \(M=[F:\mathbb Q]\).
Retain its integer \(\Delta\) and its two conductors
\(f_3,f_{-2}\) associated to (FQ.4).
These data are fixed before the new prime below.

Set
\[
 B_0=\operatorname{lcm}(24,4M,|Q|,\Delta,f_3,f_{-2}),
 \quad p\ge7,\quad p\nmid B_0,\quad
 N_0=pB_0,
 \tag{AQ.2}
\]
where \(p\) is any prime satisfying the stated
condition. Take
\(C=\operatorname{lcm}(24,4M,|Q|,p)\).
Then \(C\mid N_0\), so these are permitted literal
parameters in A.72. Its exponents are
\[
 n_j=N_0(j+4)!,\qquad
 v_p(n_0)=1.
 \tag{AQ.3}
\]
Initially its old component contains the actual
rational point \((y_0,1,1)\).
Choose that point at every old local place.
Thus the initial CRT roots are literally
\(\xi_v=\eta_v=1\).
If \(P\) is the first CRT modulus of (FQ.12),
the required precision at \(p\) is at least
\(2v_p(N_0\,6!)+4=6\).
The original first-prime congruence is consequently
\[
 q\equiv1\pmod P,\qquad p^2\mid P.
 \tag{AQ.4}
\]
No prior local or component condition has been changed.

Let \(L_p\) be the point-three specialization field
of the tangentially marked class-two exponent-\(p\)
Heisenberg quotient. Over \(\mathbb Q(\mu_p)\)
one can specify it by
\[
 x^p=3,\quad y^p=-2,\quad
 z^p=\prod_{i=0}^{p-1}(1-\zeta^i x)^i.
 \tag{AQ.5}
\]
The Kummer identity and the two noncommuting
automorphisms used for \(p=5\) in
Theorem~\ref{thm:ghost-first-integral-class3}
hold with the same formulas for every odd \(p\):
under \(x\mapsto\zeta x\), the product on the
right is multiplied by \((1-x)^p/(-2)\).
The two independent Kummer classes at the primes
three and two give the full abelian quotient
\(\mathbb F_p^2\); their commutator gives the
full center. Here the central radical really has
degree \(p\). If its displayed radicand were
\(w^p\) in \(\mathbb Q(\mu_p,x,y)\), the commuting
automorphisms \(\sigma(x)=\zeta x\),
\(\tau(y)=\zeta y\) of that abelian field would obey
\(\sigma(w)=\zeta^a w(1-x)/y\) and
\(\tau(w)=\zeta^b w\) for some integers \(a,b\).
Their two composites on \(w\) differ by \(\zeta\),
a contradiction. The identical calculation over
\(\mathbb Q(\mu_p,t)\) proves connectedness of the
generic Heisenberg cover. Thus
\[
 \operatorname{Gal}(L_p/\mathbb Q)
       =H_p\rtimes\mathbb F_p^\times,\qquad
 (L_p)^{\rm ab}=\mathbb Q(\mu_p).
 \tag{AQ.6}
\]
The action on the two abelian coordinates is scalar
cyclotomic and on the center is its square.
This is the actual point-three quotient section:
the generic equations with \(3,-2\) replaced by
\(t,1-t\) have the tangential branch \(y(0)=z(0)=1\).
Write \(\Theta(X)=\prod_{i=0}^{p-1}(1-\zeta^iX)^i\).
For \(\lambda\in\mathbb F_p^\times\), let \(k\)
represent \(\lambda^{-1}\), and let \(i_j\) be the
representative of \(\lambda^{-1}j\) in
\(\{0,\ldots,p-1\}\). Reindexing the product gives
\[
 \Theta(X)\big|_{\zeta\mapsto\zeta^\lambda}
   =\Theta(X)^k
       \left(\prod_{j=0}^{p-1}
          (1-\zeta^jX)^{(i_j-kj)/p}\right)^p.
 \tag{AQ.6a}
\]
The exponents in the parenthesis are integers.
The resulting semilinear automorphism fixes \(X,Y\)
and sends \(Z\) to \(Z^k\) times that parenthesis.
It preserves the tangential constant term one.
Such automorphisms compose exactly: their possible
discrepancy is a central \(p\)-th root of unity,
which this constant term forces to be one.
Conjugation multiplies both inertia generators
by \(\lambda\), and hence their commutator by
\(\lambda^2\). Thus the generic cover and its
specialization have the asserted actual marking.
We shall also use the cohomology vanishing
\(H^1(G_{\mathbb Q,\{p,\infty\}},\mathbb F_p(2))=0\),
whose calculation for all \(p\ge5\) is given in
the proof of Lemma~\ref{lem:ghost-finite-prefix-determines}.

Write
\[
 K=F\mathbb Q(\mu_P),\qquad
 R_b=K((-2)^{1/n_0}),\qquad
 R'=R_b(3^{1/p}).
 \tag{AQ.7}
\]
By (AQ.6), \(L_p\cap\mathbb Q(\mu_P)=\mathbb Q(\mu_p)\).
The extension \(K/\mathbb Q(\mu_P)\) has degree
dividing \(M\), prime to \(p\).
Since \(L_p\mathbb Q(\mu_P)/\mathbb Q(\mu_P)\)
is a \(p\)-extension, the two extensions are
linearly disjoint. In particular
\(\operatorname{Gal}(L_pK/K)=H_p\) and
\(L_p\cap K=\mathbb Q(\mu_p)\).

The classes of three and minus two are independent
modulo \(p\)-th powers in \(\mathbb Q(\mu_P)\):
a nontrivial rational combination has a prime
valuation not divisible by \(p\), hence its real
pure radical has degree \(p\) and is not normal
over \(\mathbb Q\). Such a radical cannot belong
to an abelian extension. They remain independent
in \(K\), by restriction followed by norm through
the extension of degree prime to \(p\).
The extension \(R_b/K\) is cyclic and has its
unique degree-\(p\) subextension
\(K((-2)^{1/p})\).
It follows that
\(R_b\cap K(3^{1/p})=K\).
Moreover \(R'/K\) is abelian and contains the
full maximal abelian subextension of \(L_pK/K\).
Therefore the exact intersection is
\[
 L_pK\cap R'=K(3^{1/p},(-2)^{1/p}).
 \tag{AQ.8}
\]

Choose a nonzero central element of \(H_p\)
on \(L_p\), the identity on \(K\), and the identity
on \(R'\).
The first two actions agree on
\(L_p\cap K=\mathbb Q(\mu_p)\).
Their combined action on \(L_pK\) agrees with
the third on (AQ.8), because a central element
fixes both radicals there.
This proves compatibility in the actual
Galois fibre product, not just degree independence.
Chebotarev \cite[Theorem 8.31]{GhostMilneANT}
supplies a fresh prime \(q\) in that conjugacy
class, exceeding every bound in A.72.
It splits in \(F\), in \(\mathbb Q(\mu_P)\),
and in \(\mathbb Q(\mu_{n_0},(-2)^{1/n_0})\),
and also splits \(3^{1/p}\).
It satisfies all the original first-prime
requirements (AQ.4) and (FQ.14).
Its point-three \(H_p\) Frobenius is central
nonzero; conjugation only multiplies that
central coordinate by a nonzero scalar.

Choose the second prime and every later prime
by (FQ.12)--(FQ.15).
The initial extra condition affects none of
its Hensel or component invariants.
The actual component morphisms (FQ.16),
the fibre-product quotient (FQ.1), and the
nested compact conjugator sets give exactly
the section and local comparisons in (AQ.1).
Its completed coordinates are
\[
 A=3\prod_{j\ge0}q_j^{n_j},\qquad
 B=-2\prod_{j\ge0}r_j^{n_j}.
 \tag{AQ.9}
\]

It remains to compute the different lifting
fibre asserted in the theorem.
Choose the literal abelian cocycles of the
convergent products (AQ.9).
Every \(n_j\) is divisible by \(p\), so they
reduce to the cocycles of \((3,-2)\) modulo
\(p\).
For any integral class-two section, its
central difference from point three is thus
an \(\mathbb F_p(2)\)-cocycle.
At a new support prime its inertia is an
integer valuation \(n_j\) times an integral
local primitive, so it is zero modulo \(p\).
At zero-valuation primes it is zero already
integrally, by the tame relation.
At the two old primes two and three, the
differences of local primitives are divisible
by \(p\): their Frobenius equations divide
the \(p\)-divisible abelian difference by
\(q'-1\), a unit modulo \(p\) for
\(q'=2,3\).
Therefore the finite central difference
is unramified outside \(p\).
The cohomology vanishing specified after
(AQ.6) makes it a global central coboundary.
Every integral class-two section consequently
has the same point-three \(H_p\) image.
At the selected \(q\), its abelian Frobenius
coordinates are zero modulo \(p\), and its
central coordinate is nonzero modulo \(p\).
These statements survive geometric conjugation
because \(q\equiv1\pmod p\) and the image
is central.

Now \(m=v_q(A)=n_0\) has \(v_p(m)=1\).
In the notation of
Proposition~\ref{prop:ghost-class3-residue},
the ordinary local unit equation makes the
second Kummer cocycle an integral local
coboundary. Its primitive \(u\) has
\(\beta=(q-1)u\), and the exact class-two
inertia calculation gives \(t=mu\) for every
integral lift.
The actual weight-three residue is
\[
 Z_{D,q}
   =m(\gamma-q\alpha u)
          +(q^2-1)u\binom m2.
 \tag{AQ.10}
\]
Here \(\alpha\equiv0\pmod p\),
\(\gamma\not\equiv0\pmod p\), and
\(v_p(q^2-1)\ge2\).
The first summand of (AQ.10) has valuation
exactly one; the second has valuation at
least three, since \(p\ge7\).
Thus \(Z_{D,q}/(q^2-1)\notin\mathbb Z_p\).
The exact strict residue criterion of that
proposition excludes every integral
class-three lift.
Corollary~\ref{cor:ghost-class2-lift} supplies
a full integral class-two section for every
ordinary pair, including (AQ.9).
Its first unrestricted integral failure is
therefore exactly class three.
\end{proof}

The program \texttt{check\_a72\_class3.py} verifies
the exact radical and descent exponent identities,
the initial factorial precision, and every residue
class in (AQ.10) modulo \(p^2\) for
\(p=7,11,13,17,19\).
It also generates all \(2058\) elements of
\(H_7\rtimes\mathbb F_7^\times\).
An independent small central-Frobenius certificate is
\[
 \begin{gathered}
 p=7,\quad q=46649,\quad \zeta=8434,\quad x=9903,\\
 \Theta(x)=20860,\qquad
 \Theta(x)^{(q-1)/7}=39280\ne1\pmod q.
 \end{gathered}
\]
All seven choices of the seventh root of three
give the same nontrivial symbol.
This small prime checks the additional Heisenberg
condition; existence with the arbitrary prescribed
A.72 boundary field and CRT data is proved by
(AQ.8), not by this small numerical certificate.
